\documentclass{report}
\usepackage{amsmath,amssymb}
\setlength{\parindent}{0mm}
\setlength{\parskip}{1em}
\begin{document}
\begin{center}
\rule{6in}{1pt} \
{\large Yasunori Aoki \\
{\bf Cluster Approximation method for Inverse Problems: applied to Model Parameter Estimation}}

University of Waterloo \\ Department of Applied Mathematics \\ 200 University Ave West \\ Waterloo \\ Ontario \\ N2L 3G1 \\ Canada
\\
{\tt yaoki@uwaterloo.ca}\\
Ken Hayami\\
Hans De Sterck\\
Ben Holder\\
Akihiko Konagaya\end{center}

We describe a new method for inverse problem that proceeds by
simultaneously finding multiple solutions of the parameter estimation
problem. Estimation of the parameter of a finitely parameterized ODE
based mathematical model from an experimental data is required in many
fields of science. This problem can be thought as solving a system of
nonlinear equations
\begin{eqnarray}
\boldsymbol{f}(\boldsymbol{\theta})&=&\boldsymbol{y}^*\,, \label{eq_1}
\end{eqnarray}
where $\boldsymbol{f}$ is a map from $\mathbb{R}^m$ to $\mathbb{R}^n$,
$\boldsymbol{\theta}$ is the parameter vector, and $\boldsymbol{y}^*$ is
the experimental data. A standard approach for solving this type of
system of nonlinear equations is a gradient based method such as the
Newton method or the Levenberg-Marquard method. However, in the case of
finding parameters from experimental data the problem may be more
complicated than just solving \eqref{eq_1} in the following ways, and
standard nonlinear solvers may not be adequate:
\begin{itemize}
\item There is no guarantee that a solution to \eqref{eq_1} exists, nor
that it is unique.
\item We often do not have a good initial guess for the parameter value
$\boldsymbol{\theta}$.
\item The experimental data contains a lot of error.
\end{itemize}
In order to overcome these challenges we have constructed a numerical
algorithm called the Cluster Approximation method that approximately
solves the system of nonlinear equations. This method differs from
traditional nonlinear equation solvers in the following ways:
\begin{itemize}
\item It find multiple possible solutions that approximately satisfy \eqref{eq_1}.
\item It is more robust against the local minima of the objective
function compare to other gradient based algorithms since it constructs a
linear approximation from multiple points.
\item It is faster than population based stochastic method as it is a
gradient base algorithm.
\end{itemize}
Simply put, the Cluster Approximation method combines the desirable
properties of a population based global optimization method (like
semi-global convergence) and with properties of gradient based method
(like convergence speed), as it moves cluster of tentitive solution
points using gradient like information.
We demonstrate these points through solving three different types of
inverse problems in mathematical biology: underdetermined inverse problem
(number of observations less than the number of parameters),
overdetermined inverse problem (number of observations more than the
number of parameters), and parameter identifiability analysis problem (an
approximation of an inverse image of a set).


\end{document}
